\newpage
\selectlanguage{greek}
\chapter*{Περίληψη}
\addcontentsline{toc}{chapter}{Περίληψη}

Αυτή η διπλωματική εργασία παρουσιάζει το \textlatin{\textbf{S\textsuperscript{4}D} (Style and Structure Similarity for Site Domains)}, μια ερευνητική πλατφόρμα για την ανάλυση ομοιοτήτων μεταξύ ιστοσελίδων βάσει της δομής και του οπτικού τους στυλ. Το σύστημα επεξεργάζεται τις σελίδες αναλύοντας τη δομή του \textlatin{DOM} και τις ιδιότητες \textlatin{CSS}, συνδυάζοντάς τα σε μία ενιαία μορφή για συγκριτική αξιολόγηση.

Οι χρήστες μπορούν να εισάγουν οποιονδήποτε αριθμό \textlatin{URLs} για ανάλυση. Κάθε σελίδα επεξεργάζεται ανεξάρτητα, με εξαγωγή δομικών και στυλιστικών πληροφοριών. Μετά το στάδιο της ομαδοποίησης (\textlatin{clustering}), η σύγκριση και αξιολόγηση ομοιότητας γίνεται βάσει του \textlatin{domain}, και όχι σε επίπεδο μεμονωμένων \textlatin{URLs}. Αυτό επιτρέπει τη μελέτη των ομοιοτήτων σε επίπεδο ιστοσελίδας συνολικά.

Για την ομαδοποίηση χρησιμοποιείται ο αλγόριθμος \textlatin{HDBSCAN}, ενώ η ποιότητα των ομάδων αξιολογείται με το \textlatin{DBCV}. Τα απομονωμένα στοιχεία επανατοποθετούνται σε κατάλληλες ομάδες όταν είναι δυνατόν. Η πλατφόρμα επιτρέπει τη διεξαγωγή πολλαπλών πειραμάτων με διαφορετικά σύνολα \textlatin{URL} και επίπεδα \textlatin{DOM}.

Το \textlatin{\textbf{S\textsuperscript{4}D}} υλοποιείται με \textlatin{React} για το περιβάλλον χρήστη, \textlatin{Flask} για τις \textlatin{API} και \textlatin{Python} για την επεξεργασία. Παράγει γραφήματα και αρχεία δεδομένων μέσω \textlatin{Gnuplot} και προσφέρει αρχείο αποτελεσμάτων για σύγκριση πειραμάτων. Εξάγει επίσης \textlatin{HTML} αρχεία με χρωματική κωδικοποίηση που αναδεικνύουν τις δομικές ομοιότητες, καθώς και πίνακες συσχετίσεων βασισμένους σε υπολογισμούς συνημιτόνου. Οι χρονικές μετρήσεις βοηθούν στην αξιολόγηση της απόδοσης του συστήματος.

Οι δοκιμές δείχνουν ότι το \textlatin{\textbf{S\textsuperscript{4}D}} ομαδοποιεί και συγκρίνει με ακρίβεια στοιχεία ιστοσελίδων βάσει δομής και εμφάνισης, προσφέροντας ένα αξιόπιστο εργαλείο για συγκριτική ανάλυση σε επίπεδο \textlatin{domain}.
